% TOP_PROBLEM: {"id": 30007029, "title": "Minimal or Mori models for generalized rational lc pairs", "category_id": 5, "category": {"id": 5, "name": "algebraic_geometry", "display_name": "Algebraic Geometry", "description": "Geometric objects defined by polynomial equations.", "slug": "algebraic-geometry", "order_index": 5, "created_at": "2026-07-31T15:26:25.671Z"}, "status": "open"}
% FIRST_POSTED: 2026-10-06
% ATTEMPT_STATUS: solved
% OPENAI_CLAIM: {"schema_version":1,"source_repo":"https://github.com/openai/math","source_commit":"adc7f1241b42e322a6451854ab7e4b4c146bf78a","release_date":"2026-10-06","match":"full","resolution":"proved","independently_reviewed":false,"second_pass":"agreed","families":[{"family":"036","title":"Numerical semiampleness and generalized minimal models","subject":"Algebraic and complex geometry","manuscripts":[{"title":"Minimal models and Mori fibre spaces for generalized log canonical ℚ-pairs","pdf_path":"preprints/Minimal-models-and-Mori-fibre-spaces-for-generalized-log-canonical-Q-pairs-September-24-2026/paper.pdf","readme_path":"preprints/Minimal-models-and-Mori-fibre-spaces-for-generalized-log-canonical-Q-pairs-September-24-2026/README.md","date":"2026-09-24","theorem_ref":"thm:main (sections/01-introduction.tex:16)","statement_sources":[{"path":"preprints/Minimal-models-and-Mori-fibre-spaces-for-generalized-log-canonical-Q-pairs-September-24-2026/build/sections/01-introduction.tex","line":16,"label":"thm:main"}]}],"lean":null,"reasoning_trace":null}]}
\documentclass[11pt]{article}
\usepackage[margin=25mm]{geometry}
\usepackage{fontspec}
\setmainfont{Latin Modern Roman}
\usepackage{amsmath,amssymb,mathtools}
\usepackage{unicode-math}
\setmathfont{Latin Modern Math}
\usepackage{enumitem}
\usepackage{xurl,hyperref,bookmark}
\hypersetup{colorlinks=true,urlcolor=blue}
\setcounter{secnumdepth}{0}
\setlength{\parindent}{0pt}
\setlength{\parskip}{5pt}
\setlength{\emergencystretch}{3em}
\begin{document}
\section*{Minimal or Mori models for generalized rational lc pairs}

\subsection{OpenAI's claim}
Projective generalized log-canonical rational pairs are claimed to admit minimal models in the pseudo-effective case and Mori fiber spaces otherwise, with the chosen nef b-divisor retained. The cited manuscript argues for that result; its relation to this problem is specified below.

\subsection{Mathematical statement of the claimed result}
\paragraph{Mathematical claim 1.} Over a field that has characteristic zero and is algebraically closed, let $(X,B+M)$ be a projective generalized log canonical rational pair: $B\ge0$ is rational, and $M$ is the trace of a rational Cartier divisor, nef on a higher birational model projective over the base, with $D=K_X+B+M$ rational Cartier. OpenAI claims a birational map $\phi:X\dashrightarrow Y$ extracting no divisor, preserving the nef b-divisor, such that $(Y,B_Y+M_Y)$ remains generalized log canonical and the discrepancy comparison is
\[
A_{X,B+M}(E)\le A_{Y,B_Y+M_Y}(E)
\]
for every prime divisor over the function field, strictly for original contracted primes. If $D$ is pseudo-effective, the endpoint adjoint is nef. Otherwise it claims a connected-fiber projective contraction $Y\to T$ with $\dim T<\dim Y$, $\rho(Y/T)=1$ and $-(K_Y+B_Y+M_Y)$ relatively ample. The inequalities use log discrepancies, equal to $1-b$ for a boundary component of coefficient $b$.

\subsection{Outline of the argument}
Manuscript 1 (family 036, 2026-09-24): The argument extends positive-boundary programs by selecting valuations with controlled normalized volume and character congruences. Equivariant regrading and finitely generated Rees constructions combine with complements and polarized boundedness to obtain the endpoint while keeping the generalized nef data fixed.

\subsection{Scope relative to this problem}
The source ranges over projective generalized log canonical rational pairs over a characteristic-zero field assumed algebraically closed, with fixed nef rational b-Cartier data. Its model extracts no divisor, has nondecreasing generalized log discrepancies with strict increase on contracted original primes, and gives the nef or relative-Picard-one Mori alternative according to pseudo-effectivity. The source field convention and rational-pair restrictions are part of this formulation.

\subsection{Formal verification}
Family 036 has no Lean scope page in this snapshot; no formalization of these manuscript results is listed.

This site has not executed the Lean files or checked their dependencies. These descriptions report the scope of the linked declarations, rather than a new verification of the proofs.

\subsection{Status caveat}
OpenAI attributes these results to an internal model. Its release notice cautions that results lacking full formalization can contain mistakes. This site has not independently assessed the proofs; a solved label records the scope of the reported claim.

\subsection{Sources}
\begin{itemize}
\item OpenAI, \emph{Minimal models and Mori fibre spaces for generalized log canonical \(\mathbb Q\)-pairs}, 2026-09-24; citation key \texttt{OAI:Minimal-models-and-Mori-fibre-spaces-for-generalized-log-canonical-Q-pairs-September-24-2026}. \href{https://github.com/openai/math/blob/adc7f1241b42e322a6451854ab7e4b4c146bf78a/preprints/Minimal-models-and-Mori-fibre-spaces-for-generalized-log-canonical-Q-pairs-September-24-2026/paper.pdf}{PDF (GitHub)}; \href{https://github.com/openai/math/raw/adc7f1241b42e322a6451854ab7e4b4c146bf78a/preprints/Minimal-models-and-Mori-fibre-spaces-for-generalized-log-canonical-Q-pairs-September-24-2026/paper.pdf}{PDF (direct)}; \href{https://github.com/openai/math/blob/adc7f1241b42e322a6451854ab7e4b4c146bf78a/preprints/Minimal-models-and-Mori-fibre-spaces-for-generalized-log-canonical-Q-pairs-September-24-2026/README.md}{Manuscript page and citation}; \href{https://github.com/openai/math/blob/main/preprints/Minimal-models-and-Mori-fibre-spaces-for-generalized-log-canonical-Q-pairs-September-24-2026/paper.pdf}{Latest version}; Statement source: \href{https://github.com/openai/math/blob/adc7f1241b42e322a6451854ab7e4b4c146bf78a/preprints/Minimal-models-and-Mori-fibre-spaces-for-generalized-log-canonical-Q-pairs-September-24-2026/build/sections/01-introduction.tex\#L16}{thm:main}.
\item \href{https://github.com/openai/math/blob/adc7f1241b42e322a6451854ab7e4b4c146bf78a/CONTENTS.md}{OpenAI catalogue entry}.
\end{itemize}
\end{document}
